% ============================================================
% 01 研究问题、任务目标与评价方式（页数预算 ≈0.6 页）
% 依据：仓库根目录 assignment.md 第 1、2、5 节及 docs/course-overview.md。
% 本节界定问题、目标与评价口径；具体技术选型与实验进展见后续章节。
% ============================================================
\section{研究问题、任务目标与评价方式}

\textbf{研究问题与任务目标。}本项目以量子一号 Pro 轮式双臂机器人为平台，研究在固定桌面场景中，如何根据“颜色\,+\,体积”指令，可靠完成指定颜色瓶子的抓取、持瓶移动、向量筒定量倾倒及收尾归位。
给定指令 $u=(c,V^\star)$，其中 $c$ 为绿、蓝、黄、红之一，$V^\star>0$ 为目标体积，机器人需从托盘 A 的四个瓶子中取出指定颜色的瓶子，移动至托盘 B，向量筒倒入目标体积的液体，并收尾归位。
任务的核心挑战在于准确识别量筒液面及其对应刻度，确定筒内液体体积，并控制倾倒过程使其接近目标值：视觉读数可能受到视角、反光与液面波动的影响，而倾倒流速随瓶身倾角和剩余液量变化，停止倾倒时还需考虑余流使最终体积超过目标值的问题。
本项目的目标是在安全完成选瓶、移动、倾倒与归位流程的基础上，提高定量倾倒的准确性，将量筒内指定颜色液体的体积相对误差控制在 $5\%$ 以内。

\textbf{评价方式。}每组进行三次测试，各次开始前公布颜色和目标体积，各组使用相同指令序列。
令第 $i$ 次的颜色、过程和体积得分分别为 $C_i,P_i,A_i$，实验报告得分为 $R\in[0,20]$，则赛道总分（满分 100 分）为
\[
  S=\frac{1}{3}\sum_{i=1}^{3}(C_i+P_i+A_i)+R.
\]
单次测试满分 80 分：量筒中液体颜色与指令一致得 $C_i=20$；过程分 $P_i$ 满分 30 分，正确抓瓶、移动至第二托盘、倒液入量筒、收尾归位分别计 7、7、9、7 分；
以测试结束后裁判读取的量筒凹液面下缘刻度为实测体积 $V_i$，相对误差为 $\varepsilon_i=|V_i-V_i^\star|/V_i^\star$，颜色正确时按下表计体积分 $A_i$（满分 30 分）：
\begin{center}
  \small
  \begin{tabular}{@{}lccccc@{}}
    \toprule
    相对误差 $\varepsilon_i$ & $[0,5\%]$ & $(5\%,10\%]$ & $(10\%,20\%]$ & $(20\%,30\%]$ & $>30\%$ \\
    \midrule
    体积分 $A_i$ & 30 & 24 & 18 & 8 & 0 \\
    \bottomrule
  \end{tabular}
\end{center}
取错颜色时 $C_i=A_i=0$，仍按完成步骤计过程分。仅量筒内液体参与计分，洒漏本身不扣分，但须及时清理。
单次不限总时长；原地卡死超过 60\,s（阈值以现场调整为准）判该次失败并终止，已完成部分仍计分。
